% Generated draft: requires mathematical and source review.
\documentclass[11pt]{article}
\usepackage[margin=25mm]{geometry}
\usepackage{fontspec}
\setmainfont{Latin Modern Roman}
\usepackage{amsmath,amssymb,mathtools,unicode-math}
\setmathfont{Latin Modern Math}
\usepackage{enumitem,needspace,xurl,hyperref}
\setlength{\parindent}{0pt}
\setlength{\parskip}{5pt}
\setlength{\emergencystretch}{3em}
\newcommand{\N}{\mathbb{N}}
\newcommand{\Z}{\mathbb{Z}}
\newcommand{\Q}{\mathbb{Q}}
\newcommand{\R}{\mathbb{R}}
\newcommand{\C}{\mathbb{C}}
\newcommand{\F}{\mathbb{F}}
\newcommand{\A}{\mathbb{A}}
\newcommand{\PP}{\mathbb{P}}
\newcommand{\E}{\mathbb{E}}
\newcommand{\eps}{\varepsilon}
\newcommand{\dd}{\,\mathrm{d}}
\begin{document}
\section*{UnsolvedMath ID: 30001301}
\textbf{Source problem code:} catalogue entry
\subsection{Definitions and mathematical statement}
For a quadratic number field \(K_D=\Q(\sqrt D)\), let \(D\) be its fundamental discriminant. Write \(C_D=\operatorname{Cl}(K_D)\) for its ordinary ideal class group, subject to the real-field convention ambiguity below. For a prime \(\ell\), define
\[
\operatorname{rk}_{\ell}(C_D)=\dim_{\F_\ell}(C_D/\ell C_D).
\]
Let \(r\in\Z_{\geq0}\). For a real number \(t>1\), put
\[
\eta_0(t)=1,\qquad \eta_k(t)=\prod_{j=1}^{k}(1-t^{-j})\quad(k\geq1),\qquad
\eta_\infty(t)=\prod_{j=1}^{\infty}(1-t^{-j}).
\]
To make one natural interpretation of the unspecified word ``density'' explicit, define
\[
\mathcal D_-(X)=\{D:\ D\text{ is a quadratic-field fundamental discriminant},\ -X\leq D<0\},
\]
\[
\mathcal D_+(X)=\{D:\ D\text{ is a quadratic-field fundamental discriminant},\ 0<D\leq X\},
\]
and, when the limits exist, define the densities conditional on the sign by
\[
\delta_{\ell,r}^{\pm}=\lim_{X\to\infty}
\frac{\#\{D\in\mathcal D_\pm(X):\operatorname{rk}_{\ell}(C_D)=r\}}
{\#\mathcal D_\pm(X)}.
\]
\textbf{Source assertion, retained literally and requiring clarification.} The supplied record asks whether, for every prime \(\ell\) and every integer \(r\geq0\), these limits exist and satisfy
\[
\delta_{\ell,r}^{-}=\ell^{-r}\eta_\infty(\ell)\eta_r(\ell)^{-2},
\qquad
\delta_{\ell,r}^{+}=\ell^{-r(r+1)}\eta_\infty(\ell)\eta_r(\ell)^{-1}\eta_{r+1}(\ell).
\]
\textbf{Unresolved qualifications.} The source does not specify the density convention; the limits above give an explicit interpretation rather than a verified source definition. It also does not distinguish the ordinary class group from the narrow class group for \(D>0\). The displayed exponents and product factors are retained exactly as supplied and require checking against the report before this can be treated as a reliable formulation of the named conjecture. The quantifier includes \(\ell=2\); no exception, conditioning, or alternative rank statistic at that prime is supplied. In particular, the background reference to four-rank results does not specify how those results relate to the stated \(\ell\)-rank assertion. No correction or missing hypothesis is inferred here.
\subsection{Short English statement}
For each prime \(\ell\) and nonnegative integer \(r\), the supplied problem proposes explicit proportions of imaginary and real quadratic-field discriminants whose class groups have \(\ell\)-rank \(r\). The formulas are reproduced literally. Their transcription, the density and real class-group conventions, and the inclusion of \(\ell=2\) require clarification. This is a statement-only draft requiring human mathematical review.
\subsection{Sources}
\begin{itemize}
\item Local UnsolvedMath registry, supplied complete source record for permanent ID \(30001301\), exact problem code \texttt{catalogue entry}, titled ``Cohen--Lenstra Densities for Class-Group \(\ell\)-Ranks,'' in the Oberwolfach Reports Open Problems collection. The proposed year is recorded as 2009; no individual proposer is supplied.
\item Supplied source citation: \emph{Explicit Methods in Number Theory (2010)}, Oberwolfach Reports, pp.~1843--1920, DOI \texttt{10.4171/owr/2009/33}. \url{https://doi.org/10.4171/owr/2009/33}. The supplied background describes these reports as records of mathematical workshops and their open-problem discussions. The citation's 2010 label and the record's 2009 year require bibliographic reconciliation.
\item Supplied literature source, labeled ``Oberwolfach report'': \url{https://ems.press/content/serial-article-files/46233?nt=1}.
\item Supplied literature source, labeled ``Fouvry--Kluners theorem'': \url{https://doi.org/10.1007/s00222-006-0021-2}.
\item Literature assessment dated exactly \texttt{2026-08-21}: the supplied curation labels the problem partially solved, cites important proved cases including quadratic-field four-rank results of Fouvry and Kluners, and says that the stated all-prime density conjecture is not proved in its full stated generality. This is supplied dated curation, not an independently verified literature assessment. The background explicitly requests independent verification of time-sensitive status and cited claims.
\item The registry labels the statement \texttt{corrected\_verified} and describes its processing as meaning-preserving semantic, typographic, and LaTeX cleanup without unsupported source-dependent additions. These supplied labels do not establish mathematical correctness or verification of the displayed formulas.
\end{itemize}
\end{document}
